\documentclass[10pt,a4paper]{amsart}
\usepackage[margin=19mm]{geometry}
\usepackage{amsmath,amssymb,mathtools}
\usepackage[scr=rsfs,cal=euler]{mathalfa}
\usepackage{xfrac}
\usepackage[unicode,hidelinks,bookmarks=false]{hyperref}
\newcommand{\ie}{i.e.\ }
\newcommand{\cf}{cf.\ }
\newcommand{\resp}{respectively\ }
\newcommand{\Subsection}{\subsection}
\providecommand{\nobreakdash}{\nobreak\hbox{-}}
\setlength{\emergencystretch}{2em}
\multlinegap0pt
\usepackage{bm}
\usepackage{bbm}
\usepackage{dsfont}
\usepackage{xspace}
\usepackage{cancel}
\usepackage{mathtools}
\usepackage{amsthm}
\makeatletter
\@ifundefined{definition}{\newtheorem{definition}{Definition}}{}
\@ifundefined{proposition}{\newtheorem{proposition}[definition]{Proposition}}{}
\makeatother
\newcommand{\R}{\mathbb{R}}
\newcommand{\dd}{\,\mathrm{d}}
\newcommand{\ii}{\mathrm{i}}
\title[Infinite-order multisoliton solutions to the Benjamin--Ono e]{Infinite-order multisoliton solutions to the Benjamin--Ono equation and soliton resolution}
\author{Louise Gassot, Patrick G\'erard}
\date{Source: arXiv:2603.15419v2 (CC BY 4.0)}
\begin{document}
\maketitle

\noindent\emph{Source and licence.} Infinite-order multisoliton solutions to the Benjamin--Ono equation and soliton resolution. Version arXiv:2603.15419v2 (CC BY 4.0), distributed under the Creative Commons Attribution 4.0 International licence, as recorded in the arXiv licence metadata for this exact version and shown on the abstract page https://arxiv.org/abs/2603.15419v2. Full rights notice: licenses/arxiv-2603.15419-CC-BY-4.0.txt.\par\smallskip

\noindent\textit{Editorial scope.} Counted unit: the Proposition whose source label is \texttt{prop:phi-domain} in the article (source0.tex lines 250--256, source label \texttt{prop:phi-domain}), delivered together with the article's own complete original proof of it (source0.tex lines 258--312, one \texttt{proof} environment of 55 lines). No mathematics is written, repaired or re-derived: statement and proof are the article's own text line for line. Required context retained uncounted: none. Every definition, display and label of those lines is retained unchanged. Omitted from this package and not redistributed: 1--249, 257--257, 313--708. Nothing inside a retained excerpt is omitted or altered: every retained body is delivered exactly as the article's lines are written, so no omission receipt of any kind accompanies this package. Citations: the retained excerpts carry no citation command at all, so no bibliography entry of the article is transcribed and no reference is redistributed. Reference adaptation: none was needed and none was made; every reference inside the delivered unit resolves inside it, so no unresolved reference or citation is produced. Printed slips kept literal: no printed word, symbol, hypothesis or number of the counted statement or of its proof was altered. Layout and class: the article's own class is \texttt{amsart} with packages including geometry, fontenc, inputenc, babel, amsmath, amssymb, amsthm, textcomp, enumitem, dsfont, mathrsfs, cite, color, hyperref; the wrapper is the lane's pinned \texttt{amsart} editorial wrapper at 10pt on A4, which restates the theorem-like environments the retained text needs and adds the article's own macro definitions that the retained text needs (\texttt{\textbackslash R}, \texttt{\textbackslash dd}, \texttt{\textbackslash ii}), with extra wrapper packages (bm, bbm, dsfont, xspace, cancel, mathtools). The delivered numbering is the unit's own. Assets: no retained span contains a figure environment, a graphics-inclusion command, a PDF-page inclusion, a table or a TikZ picture; no image file is copied into this package, the package declares no asset dependency and no image file is redistributed with it. Licence: version arXiv:2603.15419v2 of this article is distributed under the Creative Commons Attribution 4.0 International licence, as recorded in the arXiv licence metadata for this exact version and shown on the abstract page https://arxiv.org/abs/2603.15419v2. A full rights notice is delivered at licenses/arxiv-2603.15419-CC-BY-4.0.txt. Characters: every retained excerpt is pure ASCII, so the delivered engine, XeLaTeX with its default Latin Modern text font, reports no missing character.
\begin{proposition}[Lax eigenfunctions belong to the domain of $X^*$]\label{prop:phi-domain}
Let $u_0\in L^2(\R)$, and let $\varphi$ be an eigenfunction of $L_{u_0}$ associated to a negative eigenvalue $\lambda<0$. Then $\varphi\in\mathcal{D}(X^*)$.
Moreover $X^*\varphi \in \mathcal {D}(L_{u_0})$ and 
\begin{equation}\label{eq:Xstarphi}
(L_{u_0}-\lambda)(X^*\varphi )=-\ii\varphi +\frac{\ii}{2\pi |\lambda|}\langle \varphi ,\Pi u_0\rangle \Pi u_0.
\end{equation}
\end{proposition}

\begin{proof}
Taking the Fourier transform of the equality $L_{u_0}\varphi=\lambda\varphi$, we find
\begin{equation}\label{eq:phi-fourier}
\forall \xi>0,\quad \widehat{\varphi}(\xi)=\frac{\widehat{u_0\varphi}(\xi)}{(\xi+|\lambda|)}.
\end{equation}
Since $u_0\varphi \in L^1$, it is straightforward to check that the right-hand side of \eqref{eq:phi-fourier} belongs to $L^2(0,\infty)$ and is continuous on $[0,\infty)$, so that
\begin{equation}\label{I+phi}
I_+(\varphi )=\widehat \varphi (0^+)=\frac{\widehat{u_0\varphi}(0)}{|\lambda|}=\frac{\langle \varphi, \Pi u_0\rangle }{|\lambda|}.
\end{equation}

In order to show that the right-hand side of \eqref{eq:phi-fourier} belongs to $H^1(0,\infty)$, we will show that the growth rate $\tau_\eta\widehat{\varphi}$, where $\tau_\eta$ is defined for $\eta>0$ as
\[
\tau_\eta \psi:=\frac{\psi(\xi+\eta)-\psi(\xi)}{\eta},
\quad \psi\in L^2(0,\infty),
\]
admits a limit as $\eta\to 0$ in $L^2(0,\infty)$.

 First notice that the right hand side of \eqref{eq:phi-fourier} also reads $A_{u_0}(\widehat \varphi) (\xi )$, where we have set
 \[ A_v \psi (\xi )=\frac{1}{2\pi (\xi +|\lambda |)}\int_0^\infty \widehat v(\xi -\zeta)\psi (\zeta)\dd\zeta ,\]
 for $v\in L^2(\R), \psi \in L^2(0,\infty )$.
Similarly to what we already observed, $A_v $ is a bounded operator from $L^2(0,\infty )$ into itself, with a norm bounded by $C\Vert v\Vert_{L^2}$, for some constant $C$ depending only on $\lambda $.

Let us fix some parameter $R>0$ and introduce the truncated function  $u_{0,R}$ defined as
\begin{equation}\label{eq:u0R}
u_{0,R}(x)=\mathbf{1}_{|x|<R}u_0(x),
\quad
\forall x\in\R.
\end{equation}
We decompose the eigenvalue equation~\eqref{eq:phi-fourier} as
\begin{equation}\label{eq:phi-fourierbis}
\forall \xi>0, \quad \widehat{\varphi}(\xi)-A_{u_0-u_{0,R}}(\widehat \varphi) (\xi)=A_{u_{0,R}}(\widehat \varphi )(\xi).
\end{equation}
Since the norm of the operator $A_{u_0-u_{0,R}}$ tends to $0$ as $R\to \infty$, we can select $R>0$ so that $\mathrm{Id}-A_{u_0-u_{0,R}}$ is invertible as an operator of $L^2(0,\infty )$. Next we identify the commutator of $\tau_\eta $ with $A_v$. A simple calculation leads to
\begin{equation}\label{eq:tauA}
[\tau_\eta, A_v]\psi (\xi)=-\frac{A_v\psi (\xi +\eta)}{\xi +|\lambda|}+\frac{1}{2\pi(\xi +|\lambda|)}\int_0^\eta \widehat v(\xi +\eta -\zeta)\psi (\zeta)\, \frac{\dd\zeta}{\eta}.
\end{equation}
If $\psi $ is continuous on $[0,\infty )$, we conclude that $[\tau_\eta ,A_v]\psi $ has a limit $\Lambda _v\psi $ in $L^2(0,\infty )$ as $\eta \to 0^+$, given by
\begin{equation}\label{eq:limit}
 \Lambda_v\psi (\xi )=-\frac{A_v\psi (\xi)}{\xi +|\lambda|}+\frac{\widehat v(\xi )\psi (0^+)}{2\pi(\xi +|\lambda|)}. 
 \end{equation}
 Coming back to \eqref{eq:phi-fourierbis}, we have
 \[
\tau_\eta\widehat{\varphi}-A_{u_0-u_{0,R}}(\tau_\eta \widehat \varphi) =\tau_\eta A_{u_{0,R}}(\widehat \varphi )+[\tau_\eta, A_{u_0-u_{0,R}}](\widehat \varphi ),\]
 and, since $xu_{0,R}\in L^2$ and $\widehat \varphi $  is continuous on $[0,\infty )$, the right hand side has a limit in $L^2(0,\infty )$ as $\eta \to 0^+$, given by
 \[-\frac{1}{\xi +|\lambda|}A_{u_0,R}(\widehat \varphi)-\ii A_{xu_{0,R}}(\widehat \varphi )+ \Lambda_{u_0-u_{0,R}}(\widehat \varphi ).\]
 Consequently,  the growth rate satisfies, as $\eta \to 0^+$, 
\[ \tau_\eta \widehat \varphi \longrightarrow (\mathrm{Id}-A_{u_0-u_{0,R}})^{-1}\left (-\frac{1}{\xi +|\lambda|}A_{u_0,R}(\widehat \varphi)-\ii A_{xu_{0,R}}(\widehat \varphi )+ \Lambda_{u_0-u_{0,R}}(\widehat \varphi )\right ).\]
This completes the proof of the first statement. For the second statement, let us apply $\tau_\eta $ directly to \eqref{eq:phi-fourier},
\[ \tau_\eta \widehat \varphi =A_{u_0}\tau_\eta \widehat \varphi +[\tau_\eta ,A_{u_0}]\widehat \varphi .\]
Taking the limit as $\eta \to 0^+$, we obtain, in view of \eqref{eq:limit} and~\eqref{eq:phi-fourier}, 
\[ \quad  \widehat \varphi '(\xi )=A_{u_0}( \widehat \varphi ')(\xi )+\Lambda_{u_0}(\widehat \varphi )(\xi)=A_{u_0}( \widehat \varphi ')(\xi )-(\xi +|\lambda|)^{-1}\widehat \varphi (\xi)+(2\pi(\xi +|\lambda|))^{-1}\widehat u_0(\xi )I_+(\varphi ).\]
This equation exactly means
\[ (L_{u_0}-\lambda)(X^*\varphi )=-\ii \varphi +\frac{\ii}{2\pi}I_+(\varphi )\Pi u_0,\]
which, in view of \eqref{I+phi}, is  \eqref{eq:Xstarphi}.
\end{proof}

\end{document}
